\begin{textsample}{POS Dim 1 – am – Score 27.00 – am/am\_ef\_52.txt}  \label{ex:f1_pos_219}
O Ensino Religioso constitui \textbf{uma} nova Área do Conhecimento com a proposta de contribuir para a superação de todo tipo de violência por motivação religiosa ( proselitismo, discriminação, preconceito... ) dentro do espaço escolar.
Neste sentido, possibilita a construção, a socialização dos conhecimentos religiosos e a reflexão sobre a realidade no esforço de compreender o ser humano na sua diversidade de crenças e convicções, oportunizando -o a conhecer e a contribuir com o processo de construção histórico, social, cultural e religioso amazônico.
O componente curricular do Ensino Religioso compreende o ser humano na sua \textbf{totalidade} e complexidade em suas dimensões biológica, psicológica, social e especialmente a espiritual.
Desconsiderar estas dimensões é fragmentá -lo, pois o ser humano não deve aceitar a realidade como algo acabado.
A religiosidade resulta da construção humana feita \textbf{historicamente} e caracteriza -se como modo de conhecimento e de compreensão do mundo exercendo grande \textbf{influência} sobre os indivíduos e a sociedade, procurando situar o ser humano na origem das coisas, nas perspectivas presentes e futuras.
É importante \textbf{salientar} \textbf{que} o conhecimento religioso, \textbf{enquanto} manifestação da \textbf{humanidade} esteja inserido no contexto escolar, estimulando a compreensão de \textbf{que} este se dá de modo \textbf{dialógico}, privilegiando reflexões sobre limites e superações nas questões ligadas à vida e no comportamento do ser humano, no sentido de orientar a sua relação ética e social.
O Ensino Religioso deve conduzir os/as estudantes pelo caminho a ser percorrido por valores humanistas, construído sobre a base \textbf{sólida} do \textbf{amor}, da fraternidade, da bondade, da honestidade, da verdade, da humildade, da justiça, da ética, do agradecimento, da confiança e, primordialmente, solidificada no respeito e na amplitude da diversidade de pensamento, comuns à todas as filosofias e crenças.
Esta ideia é confirmada por Castella : > “ \textbf{Uma} das \textbf{tarefas} da Escola quanto ao componente curricular do Ensino Religioso é fornecer \textbf{um} instrumento de leitura da realidade e criar condições para melhorar a convivência entre as pessoas pelo conhecimento, isto é, construir os pressupostos para o diálogo ” ( 2004, p.101 ). \textbf{Historicamente}, para \textbf{que} os conhecimentos religiosos fossem reconhecidos como parte integrante da vida humana e para \textbf{que} o Ensino Religioso chegasse a \textbf{categoria} de Área do Conhecimento, foi essencial a ação de \textbf{inúmeros} \textbf{atores} notáveis e desconhecidos \textbf{que} trabalharam arduamente para lhe dar credibilidade e qualidade.
Com o processo da Constituinte de 1988, o Ensino Religioso foi efetivado como Componente curricular do Ensino Fundamental e deu continuidade à sua construção como Área do Conhecimento, a partir da escola e não através de \textbf{uma} crença ou convicções.
O Ensino Religioso Escolar tem entre seus princípios a Declaração universal dos Direitos Humanos \textbf{que} considera no Art. 18º, os diversos fatores imersos no pluralismo religioso da sociedade \textbf{contemporânea} : > “ Toda a pessoa tem direito à liberdade de pensamento, de consciência e de religião ; este direito \textbf{implica} a liberdade de mudar de religião ou de convicção, assim como a liberdade de manifestar a religião ou convicção, sozinho ou em comum, tanto em público como em privado, pelo ensino, pela prática, pelo culto e pelos ritos ” ( DUDH, 1948 ).
Neste sentido, o Ensino Religioso bem como seus marcos legais, sustentam -se em realidades \textbf{que} vão além dos fenômenos religiosos aparentes para se \textbf{ancorarem} nos costumes, hábitos e tradições da sociedade, considerando a diversidade étnicorracial, realidades estas reconhecidas na Constituição de 1988 e na LDB de 1996.
Neste sentido, Cândido afirma : > O Ensino Religioso, no contexto da educação cidadã, tem como pressuposto a \textbf{dignidade} da pessoa humana, independente da opção religiosa.
Sua especificidade consiste em trabalhar o"Fenômeno Religioso".
O reconhecimento das diferentes tradições religiosas, bem como o estudo das diferentes tradições, precisa acontecer já nas séries iniciais ( CÂNDIDO, 2002, p.40 ).
A religião deve ser entendida como \textbf{um} fenômeno \textbf{que} tem autonomia e liberdade em todos os aspectos fenomenológico, histórico, sociológico, psicológico, \textbf{antropológico} e linguístico, isto é, o fato religioso sobrevive por si só, embora esta existência esteja intimamente conectada com outros aspectos e dimensões da vida de cada pessoa e da existência coletiva da \textbf{humanidade} e no contexto amazônico.
No Brasil o preconceito e a discriminação, praticados por meio de injúria, são considerados crimes e estão previstos pelo Código Penal.
Mesmo tipificados no Código Penal como crime, convive -se com formas \textbf{explícitas} e veladas de preconceitos, discriminações e racismos, dentre estes o religioso, \textbf{uma} realidade presente também no Estado do Amazonas.
Trabalhar pelo fim das desigualdades e repensar as relações tendo como base a \textbf{equidade} é fundamental para reconhecer a diversidade religiosa como \textbf{uma} de suas maiores riquezas, tendo como pressuposto a ação pedagógica interdisciplinar.
O Ensino Religioso Escolar no Referencial Curricular do Estado do Amazonas é entendido como direito do/a estudante e com a obrigatoriedade de oferta por parte do Estado \textbf{que} está intrinsecamente ligado aos valores da democracia, da paz, dos direitos civis e políticos de cada cidadão, bem como dos Direitos Humanos.
O Brasil é \textbf{um} país \textbf{que} não possui \textbf{uma} religião oficial, mas se esforça por garantir a todos os seus cidadãos a liberdade de professar ou não \textbf{um} credo religioso como afirma na Constituição, \textbf{que} “ Garante a inviolabilidade de \textbf{uma} liberdade de consciência e crença, assegurando o livre Exercício dos cultos religiosos e a proteção aos locais de culto e suas liturgias ” ( Artigo 5º, Inciso VI ).
A proposta de trabalhar a construção dos conhecimentos religiosos com os/as estudantes na perspectiva do reconhecimento, do respeito e da valorização da diversidade religiosa nada mais é do \textbf{que} pôr em pratica o \textbf{que} está estabelecido no art. 5º da Constituição Federal e na Lei de Diretrizes e Bases da Educação ( Lei Nº 9394/96 ) : > Favorecer o desenvolvimento da capacidade de aprendizagem, tendo em vista a aquisição de conhecimentos e habilidades e a formação de atitudes e valores ; assim como o fortalecimento dos vínculos de família, dos laços de solidariedade humana e de tolerância recíproca em \textbf{que} se assenta a vida social ( LDB, 1996, p.05 ).
Neste sentido, as práticas educativas sem conhecimento científico, vazios de significado, desconectados do novo contexto legal e sociocultural, construídas durante \textbf{séculos}, precisam ser ressignificadas no sentido de \textbf{fomentar} a construção de conhecimentos no \textbf{que} tange aos diversos aspectos do Fenômeno Religioso, de forma a contribuir para \textbf{que} a atuação docente promova transformações e contribua para \textbf{que} a escola se torne \textbf{um} bem público, pertencente a todos, capaz de colaborar na promoção e valorização da diversidade religiosa.
Nesse contexto, o Ensino Religioso contribui com o fazer pedagógico nas várias ações relacionadas a essa temática, valorizando como se deu a formação do povo brasileiro, a diversidade religiosa amazônidas, sua história e cultura para construção de \textbf{uma} sociedade pluricultural e democrática.
Segundo Giliz ( 2009, p. 42 ), “ O Ensino Religioso não seria mais abordado como tema transversal, mas como área do conhecimento e, portanto, elemento indispensável ao pleno exercício da cidadania ”.
Assim sendo, o Referencial Curricular do Estado do Amazonas dá ênfase para o olhar científico sobre as religiões e as múltiplas \textbf{faces} das convicções filosóficas, vislumbrando a construção de práticas pedagógicas para o Ensino Religioso, como possibilidade de contribuir para atitudes de respeito, compreendendo o mosaico das religiosidades dos povos brasileiros/amazônicos.

% matched lemmas: amor, ancorar, antropológico, ator, categoria, contemporâneo, dialógico, dignidade, enquanto, equidade, explícito, face, fomentar, historicamente, humanidade, implicar, influência, inúmero, que, salientar, século, sólido, tarefa, totalidade, um
\end{textsample}
